\section{Staircase graphs}
\label{sec:staircase}

This section constructs, for every $n\ge2$, an ordered graph $\Gam_n$
with $s=Q_n$ on which Algorithm~H performs $n$ rounds, and proves exactly
what the algorithm does on it. The graph of Fig.~\ref{fig:gamma5} is
$\Gam_5$. The idea, already visible in Section~\ref{sec:worked}, is to let
the basic paths have different levels, and to arrange the input so that
the depth-first search always finds one SAP that crosses every shortest
basic path. That single path blocks all the others.

\subsection{The graph and its input}
\label{sec:gamma}

Fix $n\ge2$. The graph $\Gam_n$ is built on a diagram of \emph{cells}
arranged in $n-1$ left-justified rows. Row~$a$ has
\[
   \ell_a=\max\Bigl(2,\ \bigfloor{\frac{n+2-a}{2}}\Bigr)\qquad(1\le a\le n-1)
\]
cells, so the numbers of cells decrease by one every two rows, except that the
last row has two cells. Equivalently, cell~$(a,c)$ (row~$a$, column~$c$)
exists if $a+2c\le n+2$, and there is one more cell~$(n-1,2)$, which we
call the \emph{corner}. For $n=5$ the rows have $3,2,2,2$ cells; for
$n=11$ they are $6,5,5,4,4,3,3,2,2,2$ (Fig.~\ref{fig:shape}).

\begin{figure}[ht]
\centering
\begin{tikzpicture}[]
\draw[fill=boyfill,draw=boydraw!70] (0.00,0.00) rectangle +(0.50,0.50);
\draw[fill=boyfill,draw=boydraw!70] (0.50,0.00) rectangle +(0.50,0.50);
\draw[fill=boyfill,draw=boydraw!70] (1.00,0.00) rectangle +(0.50,0.50);
\draw[fill=boyfill,draw=boydraw!70] (1.50,0.00) rectangle +(0.50,0.50);
\draw[fill=boyfill,draw=boydraw!70] (2.00,0.00) rectangle +(0.50,0.50);
\draw[fill=boyfill,draw=boydraw!70] (2.50,0.00) rectangle +(0.50,0.50);
\node[font=\scriptsize,anchor=east] at (-.05,0.25) {6};
\draw[fill=boyfill,draw=boydraw!70] (0.00,-0.50) rectangle +(0.50,0.50);
\draw[fill=boyfill,draw=boydraw!70] (0.50,-0.50) rectangle +(0.50,0.50);
\draw[fill=boyfill,draw=boydraw!70] (1.00,-0.50) rectangle +(0.50,0.50);
\draw[fill=boyfill,draw=boydraw!70] (1.50,-0.50) rectangle +(0.50,0.50);
\draw[fill=boyfill,draw=boydraw!70] (2.00,-0.50) rectangle +(0.50,0.50);
\node[font=\scriptsize,anchor=east] at (-.05,-0.25) {5};
\draw[fill=boyfill,draw=boydraw!70] (0.00,-1.00) rectangle +(0.50,0.50);
\draw[fill=boyfill,draw=boydraw!70] (0.50,-1.00) rectangle +(0.50,0.50);
\draw[fill=boyfill,draw=boydraw!70] (1.00,-1.00) rectangle +(0.50,0.50);
\draw[fill=boyfill,draw=boydraw!70] (1.50,-1.00) rectangle +(0.50,0.50);
\draw[fill=boyfill,draw=boydraw!70] (2.00,-1.00) rectangle +(0.50,0.50);
\node[font=\scriptsize,anchor=east] at (-.05,-0.75) {5};
\draw[fill=boyfill,draw=boydraw!70] (0.00,-1.50) rectangle +(0.50,0.50);
\draw[fill=boyfill,draw=boydraw!70] (0.50,-1.50) rectangle +(0.50,0.50);
\draw[fill=boyfill,draw=boydraw!70] (1.00,-1.50) rectangle +(0.50,0.50);
\draw[fill=boyfill,draw=boydraw!70] (1.50,-1.50) rectangle +(0.50,0.50);
\node[font=\scriptsize,anchor=east] at (-.05,-1.25) {4};
\draw[fill=boyfill,draw=boydraw!70] (0.00,-2.00) rectangle +(0.50,0.50);
\draw[fill=boyfill,draw=boydraw!70] (0.50,-2.00) rectangle +(0.50,0.50);
\draw[fill=boyfill,draw=boydraw!70] (1.00,-2.00) rectangle +(0.50,0.50);
\draw[fill=boyfill,draw=boydraw!70] (1.50,-2.00) rectangle +(0.50,0.50);
\node[font=\scriptsize,anchor=east] at (-.05,-1.75) {4};
\draw[fill=boyfill,draw=boydraw!70] (0.00,-2.50) rectangle +(0.50,0.50);
\draw[fill=boyfill,draw=boydraw!70] (0.50,-2.50) rectangle +(0.50,0.50);
\draw[fill=boyfill,draw=boydraw!70] (1.00,-2.50) rectangle +(0.50,0.50);
\node[font=\scriptsize,anchor=east] at (-.05,-2.25) {3};
\draw[fill=boyfill,draw=boydraw!70] (0.00,-3.00) rectangle +(0.50,0.50);
\draw[fill=boyfill,draw=boydraw!70] (0.50,-3.00) rectangle +(0.50,0.50);
\draw[fill=boyfill,draw=boydraw!70] (1.00,-3.00) rectangle +(0.50,0.50);
\node[font=\scriptsize,anchor=east] at (-.05,-2.75) {3};
\draw[fill=boyfill,draw=boydraw!70] (0.00,-3.50) rectangle +(0.50,0.50);
\draw[fill=boyfill,draw=boydraw!70] (0.50,-3.50) rectangle +(0.50,0.50);
\node[font=\scriptsize,anchor=east] at (-.05,-3.25) {2};
\draw[fill=boyfill,draw=boydraw!70] (0.00,-4.00) rectangle +(0.50,0.50);
\draw[fill=boyfill,draw=boydraw!70] (0.50,-4.00) rectangle +(0.50,0.50);
\node[font=\scriptsize,anchor=east] at (-.05,-3.75) {2};
\draw[fill=boyfill,draw=boydraw!70] (0.00,-4.50) rectangle +(0.50,0.50);
\draw[fill=sapcol!45,draw=boydraw!70] (0.50,-4.50) rectangle +(0.50,0.50);
\node[font=\scriptsize,anchor=east] at (-.05,-4.25) {2};
\node[font=\scriptsize,anchor=west] at (1.10,-4.25) {$\leftarrow$ corner};
\node[font=\scriptsize,anchor=north] at (0.75,-4.60) {$n=11$, $s=36$};
\end{tikzpicture}

\caption{The cell diagram of $\Gam_{11}$; the numbers on the left are the
numbers of cells in the rows. The shaded cell is the corner cell~$(10,2)$.}
\label{fig:shape}
\end{figure}

Each cell~$(a,c)$ holds a girl~$g_{a,c}$ and a boy~$b_{a,c}$. The edges are
\begin{itemize}
\item $g_{a,c}\sg b_{a,c}$ for every cell (the \emph{inside} edges);
\item $b_{a,c}\sg g_{a,c+1}$ whenever both cells exist (\emph{horizontal}
  edges, from a boy to the girl of the next cell in his row);
\item $b_{a,c}\sg g_{a+1,c}$ whenever both cells exist (\emph{vertical}
  edges, from a boy to the girl of the cell below), \emph{except} the
  single edge $b_{n-2,2}\sg g_{n-1,2}$ into the corner, which is omitted.
\end{itemize}
The omission is a convenience for the proofs. With that edge present,
Algorithm~H still performs $n$ rounds of one SAP each (we have checked
this for $n\le30$), but in every round with at least two free girls the
corner girl would join the list $\A{BLINK}$ of a free boy, which breaks
the clean description in Lemma~\ref{lem:bfs} and changes the counts of
Theorem~W.
So a girl is adjacent to at most three boys: the boy to her left, the boy
above her, and her own. In Fig.~\ref{fig:gamma5} girls sit at the
upper left of their cells and boys at the lower right, so horizontal
edges run up and to the right, and vertical edges down and to the left.

\paragraph{The input order.}
Number the boys~$1,2,\ldots,s$ in row-major order of their cells (row by
row, left to right). Number the first girl of each row, $g_{a,1}$, as
$n-a$; number the remaining girls~$n,n+1,\ldots,s$ in row-major order.
Each girl's list consists of
\[
  \text{her left boy } b_{a,c-1},\quad\text{then her upper boy } b_{a-1,c},
  \quad\text{then her own boy } b_{a,c},
\]
omitting those that do not exist or are not adjacent to her. For $n=5$
this gives the lists~\eqref{eq:gamma5lists}.

\begin{proposition}
\label{prop:gamma}
The graph $\Gam_n$ is connected and planar, and has maximum degree three.
Its inside edges form a perfect matching, and
\[
   M=N=s=Q_n,\qquad E=3Q_n-n-\bigfloor{\frac{n+1}2}.
\]
Step~\Hs1 couples every horizontal edge $b_{a,c}\db g_{a,c+1}$. It leaves
free exactly the first girl and the last boy of every row, and puts the
free girls into \A{QUEUE} in row order~$g_{1,1},g_{2,1},\ldots,g_{n-1,1}$.
It examines $s+n-2$ adjacency entries.
\end{proposition}

\begin{proof}
Substituting $j=n+2-a$, the numbers of cells in the rows are $\max(2,\floor{j/2})$ for
$3\le j\le n+1$. Since $\sum_{j=1}^{m}\floor{j/2}=\floor{m^2/4}$, their sum is
$\floor{(n+1)^2/4}-\floor{1/2}-\floor{2/2}+1=Q_n$. There are $s$ inside
edges, $s-(n-1)$ horizontal edges, and one vertical edge into every cell
outside row~1 except the corner, so $s-\ell_1-1$ vertical edges when
$n\ge3$; with $\ell_1=\floor{(n+1)/2}$ this gives $E$ (and $n=2$ is
checked directly). Girls have at most three neighbors by construction,
and so do boys (their own girl, the girl to the right, and the girl
below). Rows are connected by their horizontal and inside edges, and the
vertical edges in column~1 connect the rows. The drawing of
Fig.~\ref{fig:gamma5} has no crossings, because every edge stays inside
its own cell or in the gap between two adjacent cells.

Step~\Hs1 scans girls~$s,s-1,\ldots,n$ first; these are the girls not in
column~1, and each takes the first boy on her list, her left boy. All
these boys are different. Then girls~$n-1,\ldots,1$, the first girls of
rows~$1,\ldots,n-1$, find their upper and own boys taken. The count of
examinations is $(s-n+1)+1+2(n-2)$.
\end{proof}

\subsection{Watching the smallest interesting case}
\label{sec:watch}

Before the general proof, look at the whole run on $\Gam_5$
(Fig.~\ref{fig:gamma5rounds}). The coloured bands are the basic paths
relative to the inside-edge perfect matching. Initially they are the rows
of the diagram, with $3,2,2,2$ cells, that is, of levels~$3,2,2,2$. The
three basic paths of level~2 are disjoint SAPs, so the dag has $L=2$ and
$\nu=3$, while $f=4$. Yet the single SAP that Algorithm~H finds runs from
the free girl of row~4 up through row~3 to the free boy of row~2, meeting
all three. After augmenting there are three basic paths, each with three
cells: the full rectangle that Section~\ref{sec:equality} demands. The
next SAP climbs a staircase through all three of them, and so on.

\begin{figure}[ht]
\centering
\resizebox{\linewidth}{!}{\begin{tikzpicture}[baseline=(current bounding box.north)]
\begin{scope}[on background layer]
\draw[hook,rowA] (0.000,0.000) -- (1.178,0.000) -- (2.356,0.000);
\draw[hook,rowB] (0.000,-1.004) -- (1.178,-1.004);
\draw[hook,rowC] (0.000,-2.009) -- (1.178,-2.009);
\draw[hook,rowD] (0.000,-3.013) -- (1.178,-3.013);
\draw[sap] (-0.223,-2.846) -- (0.223,-2.176) -- (0.955,-1.841) -- (1.401,-1.172);
\end{scope}
\draw[edge] (-0.223,-2.846) -- (0.223,-2.176);
\draw[pedge] (-0.223,-2.846) -- (0.223,-3.181);
\draw[edge] (-0.223,-1.841) -- (0.223,-1.172);
\draw[pedge] (-0.223,-1.841) -- (0.223,-2.176);
\draw[edge] (-0.223,-0.837) -- (0.223,-0.167);
\draw[pedge] (-0.223,-0.837) -- (0.223,-1.172);
\draw[pedge] (-0.223,0.167) -- (0.223,-0.167);
\draw[couple] (0.955,0.167) -- (0.223,-0.167);
\draw[pedge] (0.955,0.167) -- (1.401,-0.167);
\draw[couple] (2.133,0.167) -- (1.401,-0.167);
\draw[pedge] (2.133,0.167) -- (2.579,-0.167);
\draw[couple] (0.955,-0.837) -- (0.223,-1.172);
\draw[edge] (0.955,-0.837) -- (1.401,-0.167);
\draw[pedge] (0.955,-0.837) -- (1.401,-1.172);
\draw[couple] (0.955,-1.841) -- (0.223,-2.176);
\draw[edge] (0.955,-1.841) -- (1.401,-1.172);
\draw[pedge] (0.955,-1.841) -- (1.401,-2.176);
\draw[couple] (0.955,-2.846) -- (0.223,-3.181);
\draw[pedge] (0.955,-2.846) -- (1.401,-3.181);
\node[small girl,free] at (-0.223,-2.846) {};
\node[small girl,free] at (-0.223,-1.841) {};
\node[small girl,free] at (-0.223,-0.837) {};
\node[small girl,free] at (-0.223,0.167) {};
\node[small girl] at (0.955,0.167) {};
\node[small girl] at (2.133,0.167) {};
\node[small girl] at (0.955,-0.837) {};
\node[small girl] at (0.955,-1.841) {};
\node[small girl] at (0.955,-2.846) {};
\node[small boy] at (0.223,-0.167) {};
\node[small boy] at (1.401,-0.167) {};
\node[small boy,free] at (2.579,-0.167) {};
\node[small boy] at (0.223,-1.172) {};
\node[small boy,free] at (1.401,-1.172) {};
\node[small boy] at (0.223,-2.176) {};
\node[small boy,free] at (1.401,-2.176) {};
\node[small boy] at (0.223,-3.181) {};
\node[small boy,free] at (1.401,-3.181) {};
\node[font=\footnotesize,anchor=north] at (1.2,-3.55) {$L=2$, $f=4$};
\node[font=\scriptsize,text=black!60,anchor=north] at (1.2,-3.98) {cells $(3,2,2,2)$};
\end{tikzpicture}
\hspace{7mm}\begin{tikzpicture}[baseline=(current bounding box.north)]
\begin{scope}[on background layer]
\draw[hook,rowA] (0.000,0.000) -- (1.178,0.000) -- (2.356,0.000);
\draw[hook,rowB] (0.000,-1.004) -- (1.178,-1.004) -- (1.178,-2.009);
\draw[hook,rowC] (0.000,-2.009) -- (0.000,-3.013) -- (1.178,-3.013);
\draw[sap] (-0.223,-1.841) -- (0.223,-1.172) -- (0.955,-0.837) -- (1.401,-0.167) -- (2.133,0.167) -- (2.579,-0.167);
\end{scope}
\draw[couple] (-0.223,-2.846) -- (0.223,-2.176);
\draw[pedge] (-0.223,-2.846) -- (0.223,-3.181);
\draw[edge] (-0.223,-1.841) -- (0.223,-1.172);
\draw[pedge] (-0.223,-1.841) -- (0.223,-2.176);
\draw[edge] (-0.223,-0.837) -- (0.223,-0.167);
\draw[pedge] (-0.223,-0.837) -- (0.223,-1.172);
\draw[pedge] (-0.223,0.167) -- (0.223,-0.167);
\draw[couple] (0.955,0.167) -- (0.223,-0.167);
\draw[pedge] (0.955,0.167) -- (1.401,-0.167);
\draw[couple] (2.133,0.167) -- (1.401,-0.167);
\draw[pedge] (2.133,0.167) -- (2.579,-0.167);
\draw[couple] (0.955,-0.837) -- (0.223,-1.172);
\draw[edge] (0.955,-0.837) -- (1.401,-0.167);
\draw[pedge] (0.955,-0.837) -- (1.401,-1.172);
\draw[edge] (0.955,-1.841) -- (0.223,-2.176);
\draw[couple] (0.955,-1.841) -- (1.401,-1.172);
\draw[pedge] (0.955,-1.841) -- (1.401,-2.176);
\draw[couple] (0.955,-2.846) -- (0.223,-3.181);
\draw[pedge] (0.955,-2.846) -- (1.401,-3.181);
\node[small girl] at (-0.223,-2.846) {};
\node[small girl,free] at (-0.223,-1.841) {};
\node[small girl,free] at (-0.223,-0.837) {};
\node[small girl,free] at (-0.223,0.167) {};
\node[small girl] at (0.955,0.167) {};
\node[small girl] at (2.133,0.167) {};
\node[small girl] at (0.955,-0.837) {};
\node[small girl] at (0.955,-1.841) {};
\node[small girl] at (0.955,-2.846) {};
\node[small boy] at (0.223,-0.167) {};
\node[small boy] at (1.401,-0.167) {};
\node[small boy,free] at (2.579,-0.167) {};
\node[small boy] at (0.223,-1.172) {};
\node[small boy] at (1.401,-1.172) {};
\node[small boy] at (0.223,-2.176) {};
\node[small boy,free] at (1.401,-2.176) {};
\node[small boy] at (0.223,-3.181) {};
\node[small boy,free] at (1.401,-3.181) {};
\node[font=\footnotesize,anchor=north] at (1.2,-3.55) {$L=3$, $f=3$};
\node[font=\scriptsize,text=black!60,anchor=north] at (1.2,-3.98) {cells $(3,3,3)$};
\end{tikzpicture}
\hspace{7mm}\begin{tikzpicture}[baseline=(current bounding box.north)]
\begin{scope}[on background layer]
\draw[hook,rowA] (0.000,0.000) -- (1.178,0.000) -- (1.178,-1.004) -- (1.178,-2.009);
\draw[hook,rowB] (0.000,-1.004) -- (0.000,-2.009) -- (0.000,-3.013) -- (1.178,-3.013);
\draw[sap] (-0.223,-0.837) -- (0.223,-0.167) -- (0.955,0.167) -- (1.401,-0.167) -- (0.955,-0.837) -- (1.401,-1.172) -- (0.955,-1.841) -- (1.401,-2.176);
\end{scope}
\draw[couple] (-0.223,-2.846) -- (0.223,-2.176);
\draw[pedge] (-0.223,-2.846) -- (0.223,-3.181);
\draw[couple] (-0.223,-1.841) -- (0.223,-1.172);
\draw[pedge] (-0.223,-1.841) -- (0.223,-2.176);
\draw[edge] (-0.223,-0.837) -- (0.223,-0.167);
\draw[pedge] (-0.223,-0.837) -- (0.223,-1.172);
\draw[pedge] (-0.223,0.167) -- (0.223,-0.167);
\draw[couple] (0.955,0.167) -- (0.223,-0.167);
\draw[pedge] (0.955,0.167) -- (1.401,-0.167);
\draw[edge] (2.133,0.167) -- (1.401,-0.167);
\draw[couple] (2.133,0.167) -- (2.579,-0.167);
\draw[edge] (0.955,-0.837) -- (0.223,-1.172);
\draw[couple] (0.955,-0.837) -- (1.401,-0.167);
\draw[pedge] (0.955,-0.837) -- (1.401,-1.172);
\draw[edge] (0.955,-1.841) -- (0.223,-2.176);
\draw[couple] (0.955,-1.841) -- (1.401,-1.172);
\draw[pedge] (0.955,-1.841) -- (1.401,-2.176);
\draw[couple] (0.955,-2.846) -- (0.223,-3.181);
\draw[pedge] (0.955,-2.846) -- (1.401,-3.181);
\node[small girl] at (-0.223,-2.846) {};
\node[small girl] at (-0.223,-1.841) {};
\node[small girl,free] at (-0.223,-0.837) {};
\node[small girl,free] at (-0.223,0.167) {};
\node[small girl] at (0.955,0.167) {};
\node[small girl] at (2.133,0.167) {};
\node[small girl] at (0.955,-0.837) {};
\node[small girl] at (0.955,-1.841) {};
\node[small girl] at (0.955,-2.846) {};
\node[small boy] at (0.223,-0.167) {};
\node[small boy] at (1.401,-0.167) {};
\node[small boy] at (2.579,-0.167) {};
\node[small boy] at (0.223,-1.172) {};
\node[small boy] at (1.401,-1.172) {};
\node[small boy] at (0.223,-2.176) {};
\node[small boy,free] at (1.401,-2.176) {};
\node[small boy] at (0.223,-3.181) {};
\node[small boy,free] at (1.401,-3.181) {};
\node[font=\footnotesize,anchor=north] at (1.2,-3.55) {$L=4$, $f=2$};
\node[font=\scriptsize,text=black!60,anchor=north] at (1.2,-3.98) {cells $(4,4)$};
\end{tikzpicture}
\hspace{7mm}\begin{tikzpicture}[baseline=(current bounding box.north)]
\begin{scope}[on background layer]
\draw[hook,rowA] (0.000,0.000) -- (0.000,-1.004) -- (0.000,-2.009) -- (0.000,-3.013) -- (1.178,-3.013);
\draw[sap] (-0.223,0.167) -- (0.223,-0.167) -- (-0.223,-0.837) -- (0.223,-1.172) -- (-0.223,-1.841) -- (0.223,-2.176) -- (-0.223,-2.846) -- (0.223,-3.181) -- (0.955,-2.846) -- (1.401,-3.181);
\end{scope}
\draw[couple] (-0.223,-2.846) -- (0.223,-2.176);
\draw[pedge] (-0.223,-2.846) -- (0.223,-3.181);
\draw[couple] (-0.223,-1.841) -- (0.223,-1.172);
\draw[pedge] (-0.223,-1.841) -- (0.223,-2.176);
\draw[couple] (-0.223,-0.837) -- (0.223,-0.167);
\draw[pedge] (-0.223,-0.837) -- (0.223,-1.172);
\draw[pedge] (-0.223,0.167) -- (0.223,-0.167);
\draw[edge] (0.955,0.167) -- (0.223,-0.167);
\draw[couple] (0.955,0.167) -- (1.401,-0.167);
\draw[edge] (2.133,0.167) -- (1.401,-0.167);
\draw[couple] (2.133,0.167) -- (2.579,-0.167);
\draw[edge] (0.955,-0.837) -- (0.223,-1.172);
\draw[edge] (0.955,-0.837) -- (1.401,-0.167);
\draw[couple] (0.955,-0.837) -- (1.401,-1.172);
\draw[edge] (0.955,-1.841) -- (0.223,-2.176);
\draw[edge] (0.955,-1.841) -- (1.401,-1.172);
\draw[couple] (0.955,-1.841) -- (1.401,-2.176);
\draw[couple] (0.955,-2.846) -- (0.223,-3.181);
\draw[pedge] (0.955,-2.846) -- (1.401,-3.181);
\node[small girl] at (-0.223,-2.846) {};
\node[small girl] at (-0.223,-1.841) {};
\node[small girl] at (-0.223,-0.837) {};
\node[small girl,free] at (-0.223,0.167) {};
\node[small girl] at (0.955,0.167) {};
\node[small girl] at (2.133,0.167) {};
\node[small girl] at (0.955,-0.837) {};
\node[small girl] at (0.955,-1.841) {};
\node[small girl] at (0.955,-2.846) {};
\node[small boy] at (0.223,-0.167) {};
\node[small boy] at (1.401,-0.167) {};
\node[small boy] at (2.579,-0.167) {};
\node[small boy] at (0.223,-1.172) {};
\node[small boy] at (1.401,-1.172) {};
\node[small boy] at (0.223,-2.176) {};
\node[small boy] at (1.401,-2.176) {};
\node[small boy] at (0.223,-3.181) {};
\node[small boy,free] at (1.401,-3.181) {};
\node[font=\footnotesize,anchor=north] at (1.2,-3.55) {$L=5$, $f=1$};
\node[font=\scriptsize,text=black!60,anchor=north] at (1.2,-3.98) {cells $(5)$};
\end{tikzpicture}
\hspace{7mm}\begin{tikzpicture}[baseline=(current bounding box.north)]
\begin{scope}[on background layer]
\end{scope}
\draw[edge] (-0.223,-2.846) -- (0.223,-2.176);
\draw[couple] (-0.223,-2.846) -- (0.223,-3.181);
\draw[edge] (-0.223,-1.841) -- (0.223,-1.172);
\draw[couple] (-0.223,-1.841) -- (0.223,-2.176);
\draw[edge] (-0.223,-0.837) -- (0.223,-0.167);
\draw[couple] (-0.223,-0.837) -- (0.223,-1.172);
\draw[couple] (-0.223,0.167) -- (0.223,-0.167);
\draw[edge] (0.955,0.167) -- (0.223,-0.167);
\draw[couple] (0.955,0.167) -- (1.401,-0.167);
\draw[edge] (2.133,0.167) -- (1.401,-0.167);
\draw[couple] (2.133,0.167) -- (2.579,-0.167);
\draw[edge] (0.955,-0.837) -- (0.223,-1.172);
\draw[edge] (0.955,-0.837) -- (1.401,-0.167);
\draw[couple] (0.955,-0.837) -- (1.401,-1.172);
\draw[edge] (0.955,-1.841) -- (0.223,-2.176);
\draw[edge] (0.955,-1.841) -- (1.401,-1.172);
\draw[couple] (0.955,-1.841) -- (1.401,-2.176);
\draw[edge] (0.955,-2.846) -- (0.223,-3.181);
\draw[couple] (0.955,-2.846) -- (1.401,-3.181);
\node[small girl] at (-0.223,-2.846) {};
\node[small girl] at (-0.223,-1.841) {};
\node[small girl] at (-0.223,-0.837) {};
\node[small girl] at (-0.223,0.167) {};
\node[small girl] at (0.955,0.167) {};
\node[small girl] at (2.133,0.167) {};
\node[small girl] at (0.955,-0.837) {};
\node[small girl] at (0.955,-1.841) {};
\node[small girl] at (0.955,-2.846) {};
\node[small boy] at (0.223,-0.167) {};
\node[small boy] at (1.401,-0.167) {};
\node[small boy] at (2.579,-0.167) {};
\node[small boy] at (0.223,-1.172) {};
\node[small boy] at (1.401,-1.172) {};
\node[small boy] at (0.223,-2.176) {};
\node[small boy] at (1.401,-2.176) {};
\node[small boy] at (0.223,-3.181) {};
\node[small boy] at (1.401,-3.181) {};
\node[font=\footnotesize,anchor=north] at (1.2,-3.55) {perfect matching};
\end{tikzpicture}
}
\caption{The run of Algorithm~H on $\Gam_5$. Each panel shows the
matching before one round, drawn as in Fig.~\ref{fig:gamma5} but without
labels; coloured bands are the current basic paths, and the orange band is
the one SAP that the round finds. The graph never changes, only the
matching. The basic paths get longer and fewer; their numbers of cells go
$(3,2,2,2)\to(3,3,3)\to(4,4)\to(5)$.}
\label{fig:gamma5rounds}
\end{figure}

\begin{center}\small
\begin{tabular}{crrrrrl}
\toprule
$L$&$f$&SAPs&$W_5$&$W_{13}$&$W_{11}$&the SAP, as couples $(g,b)$ in path order\\
\midrule
2&4&1&17&5&8&$(1,6),(8,5)$\\
3&3&1&19&8&9&$(2,4),(7,2),(6,3)$\\
4&2&1&17&8&8&$(3,1),(5,2),(7,5),(8,7)$\\
5&1&1&9&5&5&$(4,1),(3,4),(2,6),(1,8),(9,9)$\\
\bottomrule
\end{tabular}
\end{center}
\noindent
The SAPs are listed by their newly formed couples, using the numerical
labels of Fig.~\ref{fig:gamma5}; for example $(1,6),(8,5)$ is
$g_1\sg b_6\db g_8\sg b_5$. Step
\Hs1 examined 12 adjacency entries, so the run costs $12+62+26+30$
elementary operations in all.

\subsection{Hooks}
\label{sec:hooks}

To describe the basic paths in general, call the state with $f$ free girls
\emph{state~$f$}; in it the dag will have level
\[
   L=n+1-f.
\]
The run goes through states~$f=n-1,n-2,\ldots,1,0$.

\begin{definition}
In state~$f$, for $1\le x\le f$, the \emph{hook} $x$ is the following
sequence of cells. Let $c=f+1-x$ be its \emph{corner column}.
\begin{itemize}
\item If row~$x$ has fewer than $c$ cells, hook~$x$ is row~$x$.
\item Otherwise hook~$x$ consists of $(x,1),(x,2),\ldots,(x,c)$, followed by
  the cells below $(x,c)$ in column~$c$, from top to bottom, except the
  corner cell.
\item The last hook, $x=f$, ends with the corner cell~$(n-1,2)$.
\end{itemize}
\end{definition}

So a hook runs east along its row to its corner column and then turns
south (Fig.~\ref{fig:hooks}). Let $M_f$ be the matching that couples
the boy of each cell of a hook to the girl of the next cell of the same
hook, and couples the girl and boy of every cell that is in no hook. A
cell in no hook is called \emph{retired}; the others are \emph{alive}.
In $M_f$ the free girls are the first girls of the hooks, the free boys
are the last boys of the hooks, and the basic paths with respect to the
inside edges are precisely the hooks. Let $M_0$ be the perfect matching
of inside edges.

\begin{figure}[ht]
\centering
\begin{tikzpicture}[]
\begin{scope}[on background layer]
\draw[hook,rowA] (0.000,0.000) -- (1.368,0.000) -- (2.736,0.000) -- (4.104,0.000) -- (5.472,0.000) -- (5.472,-1.166) -- (5.472,-2.333);
\draw[hook,rowB] (0.000,-1.166) -- (1.368,-1.166) -- (2.736,-1.166) -- (4.104,-1.166) -- (4.104,-2.333) -- (4.104,-3.499) -- (4.104,-4.666);
\draw[hook,rowC] (0.000,-2.333) -- (1.368,-2.333) -- (2.736,-2.333) -- (2.736,-3.499) -- (2.736,-4.666) -- (2.736,-5.832) -- (2.736,-6.998);
\draw[hook,rowD] (0.000,-3.499) -- (1.368,-3.499) -- (1.368,-4.666) -- (1.368,-5.832) -- (1.368,-6.998) -- (1.368,-8.165) -- (1.368,-9.331);
\draw[hook,rowE] (0.000,-4.666) -- (0.000,-5.832) -- (0.000,-6.998) -- (0.000,-8.165) -- (0.000,-9.331) -- (0.000,-10.498) -- (1.368,-10.498);
\draw[sap] (-0.259,-4.471) -- (0.259,-3.694) -- (1.109,-3.305) -- (1.627,-2.527) -- (2.477,-2.138) -- (2.995,-1.361) -- (3.845,-0.972) -- (4.363,-0.194) -- (5.213,0.194) -- (5.731,-0.194) -- (5.213,-0.972) -- (5.731,-1.361) -- (5.213,-2.138) -- (5.731,-2.527);
\end{scope}
\draw[couple] (-0.259,-10.303) -- (0.259,-9.526);
\draw[pedge] (-0.259,-10.303) -- (0.259,-10.692);
\draw[couple] (-0.259,-9.137) -- (0.259,-8.359);
\draw[pedge] (-0.259,-9.137) -- (0.259,-9.526);
\draw[couple] (-0.259,-7.970) -- (0.259,-7.193);
\draw[pedge] (-0.259,-7.970) -- (0.259,-8.359);
\draw[couple] (-0.259,-6.804) -- (0.259,-6.026);
\draw[pedge] (-0.259,-6.804) -- (0.259,-7.193);
\draw[couple] (-0.259,-5.638) -- (0.259,-4.860);
\draw[pedge] (-0.259,-5.638) -- (0.259,-6.026);
\draw[edge] (-0.259,-4.471) -- (0.259,-3.694);
\draw[pedge] (-0.259,-4.471) -- (0.259,-4.860);
\draw[edge] (-0.259,-3.305) -- (0.259,-2.527);
\draw[pedge] (-0.259,-3.305) -- (0.259,-3.694);
\draw[edge] (-0.259,-2.138) -- (0.259,-1.361);
\draw[pedge] (-0.259,-2.138) -- (0.259,-2.527);
\draw[edge] (-0.259,-0.972) -- (0.259,-0.194);
\draw[pedge] (-0.259,-0.972) -- (0.259,-1.361);
\draw[pedge] (-0.259,0.194) -- (0.259,-0.194);
\draw[couple] (1.109,0.194) -- (0.259,-0.194);
\draw[pedge] (1.109,0.194) -- (1.627,-0.194);
\draw[couple] (2.477,0.194) -- (1.627,-0.194);
\draw[pedge] (2.477,0.194) -- (2.995,-0.194);
\draw[couple] (3.845,0.194) -- (2.995,-0.194);
\draw[pedge] (3.845,0.194) -- (4.363,-0.194);
\draw[couple] (5.213,0.194) -- (4.363,-0.194);
\draw[pedge] (5.213,0.194) -- (5.731,-0.194);
\draw[edge] (6.581,0.194) -- (5.731,-0.194);
\draw[couple] (6.581,0.194) -- (7.099,-0.194);
\draw[couple] (1.109,-0.972) -- (0.259,-1.361);
\draw[edge] (1.109,-0.972) -- (1.627,-0.194);
\draw[pedge] (1.109,-0.972) -- (1.627,-1.361);
\draw[couple] (2.477,-0.972) -- (1.627,-1.361);
\draw[edge] (2.477,-0.972) -- (2.995,-0.194);
\draw[pedge] (2.477,-0.972) -- (2.995,-1.361);
\draw[couple] (3.845,-0.972) -- (2.995,-1.361);
\draw[edge] (3.845,-0.972) -- (4.363,-0.194);
\draw[pedge] (3.845,-0.972) -- (4.363,-1.361);
\draw[edge] (5.213,-0.972) -- (4.363,-1.361);
\draw[couple] (5.213,-0.972) -- (5.731,-0.194);
\draw[pedge] (5.213,-0.972) -- (5.731,-1.361);
\draw[couple] (1.109,-2.138) -- (0.259,-2.527);
\draw[edge] (1.109,-2.138) -- (1.627,-1.361);
\draw[pedge] (1.109,-2.138) -- (1.627,-2.527);
\draw[couple] (2.477,-2.138) -- (1.627,-2.527);
\draw[edge] (2.477,-2.138) -- (2.995,-1.361);
\draw[pedge] (2.477,-2.138) -- (2.995,-2.527);
\draw[edge] (3.845,-2.138) -- (2.995,-2.527);
\draw[couple] (3.845,-2.138) -- (4.363,-1.361);
\draw[pedge] (3.845,-2.138) -- (4.363,-2.527);
\draw[edge] (5.213,-2.138) -- (4.363,-2.527);
\draw[couple] (5.213,-2.138) -- (5.731,-1.361);
\draw[pedge] (5.213,-2.138) -- (5.731,-2.527);
\draw[couple] (1.109,-3.305) -- (0.259,-3.694);
\draw[edge] (1.109,-3.305) -- (1.627,-2.527);
\draw[pedge] (1.109,-3.305) -- (1.627,-3.694);
\draw[edge] (2.477,-3.305) -- (1.627,-3.694);
\draw[couple] (2.477,-3.305) -- (2.995,-2.527);
\draw[pedge] (2.477,-3.305) -- (2.995,-3.694);
\draw[edge] (3.845,-3.305) -- (2.995,-3.694);
\draw[couple] (3.845,-3.305) -- (4.363,-2.527);
\draw[pedge] (3.845,-3.305) -- (4.363,-3.694);
\draw[edge] (1.109,-4.471) -- (0.259,-4.860);
\draw[couple] (1.109,-4.471) -- (1.627,-3.694);
\draw[pedge] (1.109,-4.471) -- (1.627,-4.860);
\draw[edge] (2.477,-4.471) -- (1.627,-4.860);
\draw[couple] (2.477,-4.471) -- (2.995,-3.694);
\draw[pedge] (2.477,-4.471) -- (2.995,-4.860);
\draw[edge] (3.845,-4.471) -- (2.995,-4.860);
\draw[couple] (3.845,-4.471) -- (4.363,-3.694);
\draw[pedge] (3.845,-4.471) -- (4.363,-4.860);
\draw[edge] (1.109,-5.638) -- (0.259,-6.026);
\draw[couple] (1.109,-5.638) -- (1.627,-4.860);
\draw[pedge] (1.109,-5.638) -- (1.627,-6.026);
\draw[edge] (2.477,-5.638) -- (1.627,-6.026);
\draw[couple] (2.477,-5.638) -- (2.995,-4.860);
\draw[pedge] (2.477,-5.638) -- (2.995,-6.026);
\draw[edge] (1.109,-6.804) -- (0.259,-7.193);
\draw[couple] (1.109,-6.804) -- (1.627,-6.026);
\draw[pedge] (1.109,-6.804) -- (1.627,-7.193);
\draw[edge] (2.477,-6.804) -- (1.627,-7.193);
\draw[couple] (2.477,-6.804) -- (2.995,-6.026);
\draw[pedge] (2.477,-6.804) -- (2.995,-7.193);
\draw[edge] (1.109,-7.970) -- (0.259,-8.359);
\draw[couple] (1.109,-7.970) -- (1.627,-7.193);
\draw[pedge] (1.109,-7.970) -- (1.627,-8.359);
\draw[edge] (1.109,-9.137) -- (0.259,-9.526);
\draw[couple] (1.109,-9.137) -- (1.627,-8.359);
\draw[pedge] (1.109,-9.137) -- (1.627,-9.526);
\draw[couple] (1.109,-10.303) -- (0.259,-10.692);
\draw[pedge] (1.109,-10.303) -- (1.627,-10.692);
\node[small girl] at (-0.259,-10.303) {};
\node[small girl] at (-0.259,-9.137) {};
\node[small girl] at (-0.259,-7.970) {};
\node[small girl] at (-0.259,-6.804) {};
\node[small girl] at (-0.259,-5.638) {};
\node[small girl,free] at (-0.259,-4.471) {};
\node[small girl,free] at (-0.259,-3.305) {};
\node[small girl,free] at (-0.259,-2.138) {};
\node[small girl,free] at (-0.259,-0.972) {};
\node[small girl,free] at (-0.259,0.194) {};
\node[small girl] at (1.109,0.194) {};
\node[small girl] at (2.477,0.194) {};
\node[small girl] at (3.845,0.194) {};
\node[small girl] at (5.213,0.194) {};
\node[small girl] at (6.581,0.194) {};
\node[small girl] at (1.109,-0.972) {};
\node[small girl] at (2.477,-0.972) {};
\node[small girl] at (3.845,-0.972) {};
\node[small girl] at (5.213,-0.972) {};
\node[small girl] at (1.109,-2.138) {};
\node[small girl] at (2.477,-2.138) {};
\node[small girl] at (3.845,-2.138) {};
\node[small girl] at (5.213,-2.138) {};
\node[small girl] at (1.109,-3.305) {};
\node[small girl] at (2.477,-3.305) {};
\node[small girl] at (3.845,-3.305) {};
\node[small girl] at (1.109,-4.471) {};
\node[small girl] at (2.477,-4.471) {};
\node[small girl] at (3.845,-4.471) {};
\node[small girl] at (1.109,-5.638) {};
\node[small girl] at (2.477,-5.638) {};
\node[small girl] at (1.109,-6.804) {};
\node[small girl] at (2.477,-6.804) {};
\node[small girl] at (1.109,-7.970) {};
\node[small girl] at (1.109,-9.137) {};
\node[small girl] at (1.109,-10.303) {};
\node[small boy] at (0.259,-0.194) {};
\node[small boy] at (1.627,-0.194) {};
\node[small boy] at (2.995,-0.194) {};
\node[small boy] at (4.363,-0.194) {};
\node[small boy] at (5.731,-0.194) {};
\node[small boy] at (7.099,-0.194) {};
\node[small boy] at (0.259,-1.361) {};
\node[small boy] at (1.627,-1.361) {};
\node[small boy] at (2.995,-1.361) {};
\node[small boy] at (4.363,-1.361) {};
\node[small boy] at (5.731,-1.361) {};
\node[small boy] at (0.259,-2.527) {};
\node[small boy] at (1.627,-2.527) {};
\node[small boy] at (2.995,-2.527) {};
\node[small boy] at (4.363,-2.527) {};
\node[small boy,free] at (5.731,-2.527) {};
\node[small boy] at (0.259,-3.694) {};
\node[small boy] at (1.627,-3.694) {};
\node[small boy] at (2.995,-3.694) {};
\node[small boy] at (4.363,-3.694) {};
\node[small boy] at (0.259,-4.860) {};
\node[small boy] at (1.627,-4.860) {};
\node[small boy] at (2.995,-4.860) {};
\node[small boy,free] at (4.363,-4.860) {};
\node[small boy] at (0.259,-6.026) {};
\node[small boy] at (1.627,-6.026) {};
\node[small boy] at (2.995,-6.026) {};
\node[small boy] at (0.259,-7.193) {};
\node[small boy] at (1.627,-7.193) {};
\node[small boy,free] at (2.995,-7.193) {};
\node[small boy] at (0.259,-8.359) {};
\node[small boy] at (1.627,-8.359) {};
\node[small boy] at (0.259,-9.526) {};
\node[small boy,free] at (1.627,-9.526) {};
\node[small boy] at (0.259,-10.692) {};
\node[small boy,free] at (1.627,-10.692) {};
\node[font=\scriptsize,anchor=east] at (-0.550,0.000) {hook 1};
\node[font=\scriptsize,anchor=east] at (-0.550,-1.166) {hook 2};
\node[font=\scriptsize,anchor=east] at (-0.550,-2.333) {hook 3};
\node[font=\scriptsize,anchor=east] at (-0.550,-3.499) {hook 4};
\node[font=\scriptsize,anchor=east] at (-0.550,-4.666) {hook 5};
\end{tikzpicture}

\caption{State~$f=5$ of $\Gam_{11}$, before the round at level~$L=7$. The
five hooks (coloured) run east along their rows and turn south at
columns~$5,4,3,2,1$. Cells in columns~$6$ and beyond are retired. All
five hooks have $7$ cells. For instance cell~$(7,2)$ lies in
hook~$\min(7,6-2)=4$, at position~$2+\max(0,7+2-6)=5$, as
Lemma~\ref{lem:hooks}(a) says. The orange staircase is the SAP that
Algorithm~H finds: it starts at the free girl of hook~5 and climbs to the
free boy of hook~1.}
\label{fig:hooks}
\end{figure}

\begin{lemma}
\label{lem:hooks}
Let $1\le f\le n-1$ and $L=n+1-f$.
\begin{enumerate}[label=(\alph*)]
\item The hooks are disjoint, and their cells are exactly the cells in
  columns~$1,\ldots,f$ together with the corner cell. An alive cell
  $(a,c)$ other than the corner lies in hook~$x=\min(a,f+1-c)$, at position
  $y=c+\max(0,a+c-f-1)$ along it; the corner is at position~$L$ of
  hook~$f$.
\item Consecutive cells of a hook are joined by a horizontal or vertical
  edge of $\Gam_n$, so $M_f$ is a matching of $\Gam_n$.
\item Hook~$x$ has $\max(\ell_x,L)$ cells. Hence hooks~$1,\ldots,k-1$
  have more than $L$ cells and hooks~$k,\ldots,f$ have exactly $L$ cells,
  where
  \[
      k=\max(1,f-L),\qquad m=f-k+1=\min(f,L+1).
  \]
\item $M_{n-1}$ is the matching produced by~\Hs1.
\end{enumerate}
\end{lemma}

\begin{proof}
(a) A cell~$(a,c)$ with $c\le f$ and $a+c\le f+1$ lies in row~$a$ at or
before the corner column~$f+1-a$ of hook~$a$, so it belongs to hook~$a$ at
position~$c$. If $a+c>f+1$, the hook whose corner column is $c$ is hook
$x=f+1-c<a$; since columns are top-justified, cell~$(x,c)$ exists, so
hook~$x$ turns south at $(x,c)$ and contains $(a,c)$ at position
$c+(a-x)$, unless $(a,c)$ is the corner. Conversely every hook stays in
columns~$\le f$ apart from the corner, and the two cases above assign each
cell to only one hook.

(b) East steps use horizontal edges and south steps use vertical edges;
the omitted edge would lead into the corner, which is entered from its
left neighbour $(n-1,1)$ instead.

(c) Write $c=f+1-x$, and note that $2\ell_x\ge n+1-x$ always. If
$\ell_x<c$, hook~$x$ is row~$x$; then $x\le f-\ell_x$, so
$2\ell_x\ge n+1-f+\ell_x$, that is, $\ell_x\ge L$. Otherwise, if $x<f$, hook
$x$ consists of $c$ cells of row~$x$ followed by the cells of column~$c$
below it. Column~$c\ge2$ ends at row~$n+2-2c$ (the corner is not counted),
so the hook has $c+(n+2-2c-x)=n+2-c-x=L$ cells. The last hook has
$1+(n-1-f)+1=L$ cells. In this second case~$\ell_x\le L$ as well: for
$x=n-1$ we have $\ell_x=2=L$, and otherwise $2\ell_x\le n+2-x$, whereas
$\ell_x\ge L+1$ and $\ell_x\ge c$ would give $2\ell_x\ge n+3-x$. Finally,
$\ell_x\ge L+1=n+2-f$ if and only if $n+2-x\ge2(n+2-f)$, that is,
$x\le f-L-1$.

(d) For $f=n-1$ the corner column of hook~$x$ is $n-x$, which exceeds
$\ell_x$ when $x\le n-3$; for the last two rows the definition also gives
the whole row. So $M_{n-1}$ couples the horizontal edges, as in
Proposition~\ref{prop:gamma}.
\end{proof}

Write $y(a,c)$ for the position of an alive cell, as given by part~(a).
Two facts about it drive the next two lemmas: $y$ never decreases when we
move from an alive cell to the alive cell below it or to its right, and it
increases by exactly one from each cell of a hook to the next. (For the
corner, $y=L$; its left neighbour has position~$L-1$, and the cell above it,
$(n-2,2)$, also has position~$L$.)

\subsection{No shortcuts}

The graph has many edges that the hooks do not use, and some of them will
be needed in later rounds. The next lemma shows that none of them
creates a shorter augmenting path now.

\begin{lemma}
\label{lem:levels}
In the matching $M_f$, the girl of an alive cell at position~$y$ of its
hook is at distance $2(y-1)$ from the free girls, and its boy at distance
$2y-1$. No retired vertex can be reached by an alternating path from a
free girl. Consequently the dag has level~$L$, and its SAPs end at the
free boys of hooks~$k,\ldots,f$.
\end{lemma}

\begin{proof}
First, retired vertices. A retired boy~$b_{a,c}$ can be entered only from
the girl to his right or the girl below him, since his own girl is his
mate. Both lie in retired cells, being in column~$c$ or $c+1$, beyond
column~$f$. (When $f=1$ the corner is alive and lies below the retired
cell~$(n-2,2)$; this is one place where the omitted edge matters.) A
retired girl can be entered only from her own boy. So no alternating path
from a free girl ever reaches a retired vertex.

Now follow an alternating path from a free girl. A step from a girl to a
boy uses an edge outside $M_f$, which goes from cell~$(a,c)$ to the same
cell, to $(a,c-1)$, or to $(a-1,c)$; by the monotonicity of positions it
does not increase the position. A step from a boy to a girl uses a couple
of $M_f$, and so increases the position by exactly one. The path starts
at position~1, so a boy at position~$y$ is reached only after at least
$y-1$ couples, that is, after at least $2y-1$ edges. The beginning of his
own hook reaches him in exactly that many. The argument for girls is the
same.

Free boys are the last boys of the hooks, at positions $\max(\ell_x,L)\ge
L$. The smallest distance to one of them is therefore $2L-1$, attained
exactly by hooks~$k,\ldots,f$.
\end{proof}

\subsection{The forced search}

Now the input order enters. Fig.~\ref{fig:straight} shows a dag of a
staircase graph with each hook drawn as a straight row; the positions
become the columns, which are exactly the levels of the dag.

\begin{figure}[ht]
\centering
\begin{tikzpicture}[]
\begin{scope}[on background layer]
\draw[sap] (-0.300,-4.680) -- (0.300,-3.900) -- (1.020,-3.900) -- (1.620,-3.120) -- (2.340,-3.120) -- (2.940,-2.340) -- (3.660,-2.340) -- (4.260,-1.560);
\fill[black!4] (4.620,0.450) rectangle (6.006,-5.130);
\end{scope}
\draw[couple] (1.620,-4.680) -- (2.340,-4.680);
\draw[couple] (0.300,-4.680) -- (1.020,-4.680);
\draw[couple] (0.300,0.000) -- (1.020,0.000);
\draw[couple] (1.620,0.000) -- (2.340,0.000);
\draw[couple] (2.940,0.000) -- (3.660,0.000);
\draw[couple] (4.260,0.000) -- (4.980,0.000);
\draw[couple] (0.300,-0.780) -- (1.020,-0.780);
\draw[couple] (1.620,-0.780) -- (2.340,-0.780);
\draw[couple] (2.940,-0.780) -- (3.660,-0.780);
\draw[couple] (4.260,-0.780) -- (4.980,-0.780);
\draw[couple] (0.300,-1.560) -- (1.020,-1.560);
\draw[couple] (1.620,-1.560) -- (2.340,-1.560);
\draw[couple] (2.940,-1.560) -- (3.660,-1.560);
\draw[couple] (0.300,-2.340) -- (1.020,-2.340);
\draw[couple] (1.620,-2.340) -- (2.340,-2.340);
\draw[couple] (2.940,-2.340) -- (3.660,-2.340);
\draw[couple] (0.300,-3.120) -- (1.020,-3.120);
\draw[couple] (1.620,-3.120) -- (2.340,-3.120);
\draw[couple] (0.300,-3.900) -- (1.020,-3.900);
\draw[couple] (2.940,-3.120) -- (3.660,-3.120);
\draw[couple] (1.620,-3.900) -- (2.340,-3.900);
\draw[couple] (2.940,-3.900) -- (3.660,-3.900);
\draw[couple] (2.940,-4.680) -- (3.660,-4.680);
\draw[arc,shorten >=2.6mm,shorten <=2.4mm] (0.300,0.000) -- (-0.300,-0.780);
\draw[arc,shorten >=2.6mm,shorten <=2.4mm] (0.300,0.000) -- (-0.300,0.000);
\draw[arc,shorten >=2.6mm,shorten <=2.4mm] (0.300,-0.780) -- (-0.300,-1.560);
\draw[arc,shorten >=2.6mm,shorten <=2.4mm] (0.300,-0.780) -- (-0.300,-0.780);
\draw[arc,shorten >=2.6mm,shorten <=2.4mm] (0.300,-1.560) -- (-0.300,-2.340);
\draw[arc,shorten >=2.6mm,shorten <=2.4mm] (0.300,-1.560) -- (-0.300,-1.560);
\draw[arc,shorten >=2.6mm,shorten <=2.4mm] (0.300,-2.340) -- (-0.300,-3.120);
\draw[arc,shorten >=2.6mm,shorten <=2.4mm] (0.300,-2.340) -- (-0.300,-2.340);
\draw[arc,shorten >=2.6mm,shorten <=2.4mm] (0.300,-3.120) -- (-0.300,-3.900);
\draw[arc,shorten >=2.6mm,shorten <=2.4mm] (0.300,-3.120) -- (-0.300,-3.120);
\draw[arc,shorten >=2.6mm,shorten <=2.4mm] (0.300,-3.900) -- (-0.300,-4.680);
\draw[arc,shorten >=2.6mm,shorten <=2.4mm] (0.300,-3.900) -- (-0.300,-3.900);
\draw[arc,shorten >=2.6mm,shorten <=2.4mm] (0.300,-4.680) -- (-0.300,-4.680);
\draw[arc,shorten >=2.6mm,shorten <=2.4mm] (1.620,0.000) -- (1.020,-0.780);
\draw[arc,shorten >=2.6mm,shorten <=2.4mm] (1.620,0.000) -- (1.020,0.000);
\draw[arc,shorten >=2.6mm,shorten <=2.4mm] (1.620,-0.780) -- (1.020,-1.560);
\draw[arc,shorten >=2.6mm,shorten <=2.4mm] (1.620,-0.780) -- (1.020,-0.780);
\draw[arc,shorten >=2.6mm,shorten <=2.4mm] (1.620,-1.560) -- (1.020,-2.340);
\draw[arc,shorten >=2.6mm,shorten <=2.4mm] (1.620,-1.560) -- (1.020,-1.560);
\draw[arc,shorten >=2.6mm,shorten <=2.4mm] (1.620,-2.340) -- (1.020,-3.120);
\draw[arc,shorten >=2.6mm,shorten <=2.4mm] (1.620,-2.340) -- (1.020,-2.340);
\draw[arc,shorten >=2.6mm,shorten <=2.4mm] (1.620,-3.120) -- (1.020,-3.900);
\draw[arc,shorten >=2.6mm,shorten <=2.4mm] (1.620,-3.120) -- (1.020,-3.120);
\draw[arc,shorten >=2.6mm,shorten <=2.4mm] (1.620,-3.900) -- (1.020,-3.900);
\draw[arc,shorten >=2.6mm,shorten <=2.4mm] (1.620,-4.680) -- (1.020,-4.680);
\draw[arc,shorten >=2.6mm,shorten <=2.4mm] (2.940,0.000) -- (2.340,-0.780);
\draw[arc,shorten >=2.6mm,shorten <=2.4mm] (2.940,0.000) -- (2.340,0.000);
\draw[arc,shorten >=2.6mm,shorten <=2.4mm] (2.940,-0.780) -- (2.340,-1.560);
\draw[arc,shorten >=2.6mm,shorten <=2.4mm] (2.940,-0.780) -- (2.340,-0.780);
\draw[arc,shorten >=2.6mm,shorten <=2.4mm] (2.940,-1.560) -- (2.340,-2.340);
\draw[arc,shorten >=2.6mm,shorten <=2.4mm] (2.940,-1.560) -- (2.340,-1.560);
\draw[arc,shorten >=2.6mm,shorten <=2.4mm] (2.940,-2.340) -- (2.340,-3.120);
\draw[arc,shorten >=2.6mm,shorten <=2.4mm] (2.940,-2.340) -- (2.340,-2.340);
\draw[arc,shorten >=2.6mm,shorten <=2.4mm] (2.940,-3.120) -- (2.340,-3.120);
\draw[arc,shorten >=2.6mm,shorten <=2.4mm] (2.940,-3.900) -- (2.340,-3.900);
\draw[arc,shorten >=2.6mm,shorten <=2.4mm] (2.940,-4.680) -- (2.340,-4.680);
\draw[arc,shorten >=2.6mm,shorten <=2.4mm] (4.260,0.000) -- (3.660,-0.780);
\draw[arc,shorten >=2.6mm,shorten <=2.4mm] (4.260,0.000) -- (3.660,0.000);
\draw[arc,shorten >=2.6mm,shorten <=2.4mm] (4.260,-0.780) -- (3.660,-1.560);
\draw[arc,shorten >=2.6mm,shorten <=2.4mm] (4.260,-0.780) -- (3.660,-0.780);
\draw[arc,shorten >=2.6mm,shorten <=2.4mm] (4.260,-1.560) -- (3.660,-2.340);
\draw[arc,shorten >=2.6mm,shorten <=2.4mm] (4.260,-1.560) -- (3.660,-1.560);
\draw[arc,shorten >=2.6mm,shorten <=2.4mm] (4.260,-2.340) -- (3.660,-2.340);
\draw[arc,shorten >=2.6mm,shorten <=2.4mm] (4.260,-3.120) -- (3.660,-3.120);
\draw[arc,shorten >=2.6mm,shorten <=2.4mm] (4.260,-3.900) -- (3.660,-3.900);
\draw[arc,shorten >=2.6mm,shorten <=2.4mm] (4.260,-4.680) -- (3.660,-4.680);
\draw[pedge,opacity=.35] (4.980,0.000) -- (5.580,0.000);
\draw[pedge,opacity=.35] (4.980,-0.780) -- (5.580,-0.780);
\node[small girl,free] at (-0.300,0.000) {};
\node[small girl] at (1.020,0.000) {};
\node[small girl] at (2.340,0.000) {};
\node[small girl] at (3.660,0.000) {};
\node[small girl] at (4.980,0.000) {};
\node[small girl,free] at (-0.300,-0.780) {};
\node[small girl] at (1.020,-0.780) {};
\node[small girl] at (2.340,-0.780) {};
\node[small girl] at (3.660,-0.780) {};
\node[small girl] at (4.980,-0.780) {};
\node[small girl,free] at (-0.300,-1.560) {};
\node[small girl] at (1.020,-1.560) {};
\node[small girl] at (2.340,-1.560) {};
\node[small girl] at (3.660,-1.560) {};
\node[small girl,free] at (-0.300,-2.340) {};
\node[small girl] at (1.020,-2.340) {};
\node[small girl] at (2.340,-2.340) {};
\node[small girl] at (3.660,-2.340) {};
\node[small girl,free] at (-0.300,-3.120) {};
\node[small girl] at (1.020,-3.120) {};
\node[small girl] at (2.340,-3.120) {};
\node[small girl] at (3.660,-3.120) {};
\node[small girl,free] at (-0.300,-3.900) {};
\node[small girl] at (1.020,-3.900) {};
\node[small girl] at (2.340,-3.900) {};
\node[small girl] at (3.660,-3.900) {};
\node[small girl,free] at (-0.300,-4.680) {};
\node[small girl] at (1.020,-4.680) {};
\node[small girl] at (2.340,-4.680) {};
\node[small girl] at (3.660,-4.680) {};
\node[small boy] at (0.300,0.000) {};
\node[small boy] at (1.620,0.000) {};
\node[small boy] at (2.940,0.000) {};
\node[small boy] at (4.260,0.000) {};
\node[small boy,free,opacity=.4] at (5.580,0.000) {};
\node[small boy] at (0.300,-0.780) {};
\node[small boy] at (1.620,-0.780) {};
\node[small boy] at (2.940,-0.780) {};
\node[small boy] at (4.260,-0.780) {};
\node[small boy,free,opacity=.4] at (5.580,-0.780) {};
\node[small boy] at (0.300,-1.560) {};
\node[small boy] at (1.620,-1.560) {};
\node[small boy] at (2.940,-1.560) {};
\node[small boy,free] at (4.260,-1.560) {};
\node[small boy] at (0.300,-2.340) {};
\node[small boy] at (1.620,-2.340) {};
\node[small boy] at (2.940,-2.340) {};
\node[small boy,free] at (4.260,-2.340) {};
\node[small boy] at (0.300,-3.120) {};
\node[small boy] at (1.620,-3.120) {};
\node[small boy] at (2.940,-3.120) {};
\node[small boy,free] at (4.260,-3.120) {};
\node[small boy] at (0.300,-3.900) {};
\node[small boy] at (1.620,-3.900) {};
\node[small boy] at (2.940,-3.900) {};
\node[small boy,free] at (4.260,-3.900) {};
\node[small boy] at (0.300,-4.680) {};
\node[small boy] at (1.620,-4.680) {};
\node[small boy] at (2.940,-4.680) {};
\node[small boy,free] at (4.260,-4.680) {};
\node[font=\scriptsize,anchor=east] at (-0.600,0.000) {hook 1};
\node[font=\scriptsize,anchor=east] at (-0.600,-0.780) {hook 2};
\node[font=\scriptsize,anchor=east] at (-0.600,-1.560) {hook 3};
\node[font=\scriptsize,anchor=east] at (-0.600,-2.340) {hook 4};
\node[font=\scriptsize,anchor=east] at (-0.600,-3.120) {hook 5};
\node[font=\scriptsize,anchor=east] at (-0.600,-3.900) {hook 6};
\node[font=\scriptsize,anchor=east] at (-0.600,-4.680) {hook 7};
\node[font=\scriptsize,text=black!60] at (0.000,0.550) {1};
\node[font=\scriptsize,text=black!60] at (1.320,0.550) {2};
\node[font=\scriptsize,text=black!60] at (2.640,0.550) {3};
\node[font=\scriptsize,text=black!60] at (3.960,0.550) {4};
\node[font=\scriptsize,text=black!60] at (5.280,0.550) {5};
\node[font=\scriptsize,text=black!60,anchor=east] at (-0.600,0.550) {position};
\end{tikzpicture}

\caption{The dag of the round at level~$L=4$ of $\Gam_{10}$ (state~$f=7$),
with the hooks straightened. Arrows are the arcs of $\A{BLINK}$, from a
boy to a girl one level down: every boy has an arc to his own girl, and
the boy at position~$y$ of hook~$x$ also has an arc to the girl below him
when $x+y\le7$; that arc comes first in $\A{BLINK}$. Hooks~1 and~2 have
more than $L$ cells (the shaded column lies beyond the dag). The search
starts at the free boy of hook~3, the first hook with exactly $L$ cells,
and always takes the first arc, so it goes down a staircase to the last
free girl.}
\label{fig:straight}
\end{figure}

\begin{lemma}
\label{lem:bfs}
Suppose a round starts with the matching $M_f$ and with the free girls in
\A{QUEUE} in hook order. Then each level~$G_j$ of the dag is listed in hook
order, and for the boy~$b$ at position~$y\le L$ of hook~$x$,
\[
  \A{BLINK}[b]=\begin{cases}
   (\text{girl at position~$y$ of hook~$x+1$},\ \text{girl at position~$y$ of hook~$x$}),&\text{if } x+y\le f;\\
   (\text{girl at position~$y$ of hook~$x$}),&\text{if } x+y>f.
  \end{cases}
\]
The roots of the depth-first search are the free boys of hooks
$k,k+1,\ldots,f$, in this order.
\end{lemma}

\begin{proof}
Let $b=b_{a,c}$ lie at position~$y$ of hook~$x$. By
Lemma~\ref{lem:levels}, the girls one level below $b$ are those at
position~$y$. Of his three possible neighbours, his own girl qualifies;
the girl to his right does not, since her position is larger; and the girl
below him, $g_{a+1,c}$, if she is adjacent to him, qualifies exactly when
the two cells have the same position. By the formula of
Lemma~\ref{lem:hooks}(a) that means $a+c\le f$, with one exception: the
corner has the same position as the cell above it, but the edge between
them is the omitted one, which is the main reason for omitting it. When
$a+c\le f$, then $(a,c)$ lies in the first part of hook~$a$, so $c=y\le L$,
and $(a+1)+2c\le f+1+L=n+2$; hence cell~$(a+1,c)$ exists, and it is not the
corner because $a+c\le f<n$. Also $(a,c)$ is at
position~$c$ of hook~$a$, so $a+c\le f$ is the condition $x+y\le f$, and
$g_{a+1,c}$ is at position~$y$ of hook~$x+1$.

Now argue by induction on $j$ that $G_j$ is in hook order. For $j=0$ this
is the hypothesis. When the girls of $G_j$ are scanned, the girl of hook
$x$ meets at most three boys: her own boy; her left boy, whose position is
smaller, so that he is already in the dag at a lower level and step~\Hs6
ignores him; and her upper boy, whose position is at most hers. (On the
vertical part of a hook her upper boy is her mate, at the level below.) The
only girls of $G_j$ adjacent to her own boy are herself and the girl of
hook~$x+1$, who comes later; so she is the one who discovers him. Thus $B_{j+1}$
is discovered in hook order, and so is $G_{j+1}$, which consists of their
mates. When the girl of hook~$x+1$ later meets the same boy as her upper
boy, step~\Hs6 inserts her at the head of his list, which gives the stated
order of $\A{BLINK}$.

At level~$L$ the boys of hooks~$1,\ldots,k-1$ are matched and are
discovered before the first free boy (in hook~$k$), so they enter the
dag; the free boys of hooks~$k,\ldots,f$ follow in order and become the
roots.
\end{proof}

\begin{lemma}
\label{lem:dfs}
Under the hypotheses of Lemma~\ref{lem:bfs}, the round finds exactly one
SAP. Its free girl is the last girl in \A{QUEUE}, and afterwards the matching
is $M_{f-1}$, with the remaining free girls still in hook order.
\end{lemma}

\begin{proof}
Write $\langle x,y\rangle$ for position~$y$ of hook~$x$, and recall that hooks
$k,\ldots,f$ have exactly $L$ cells and $f-k\le L$.

\emph{The first root.} The search starts at the boy~$\langle k,L\rangle$ and always
follows the first remaining arc. Suppose first that $f\ge L+1$, so that
$k=f-L$ and $k+L=f$. Then the first arc leads to the girl~$\langle k+1,L\rangle$; her
mate is the boy~$\langle k+1,L-1\rangle$, whose hook and position again add up to $f$,
so the next arc leads to the girl~$\langle k+2,L-1\rangle$, and so on down the
staircase
\[
   \text{boy }\langle i,f-i\rangle\ \to\ \text{girl }\langle i+1,f-i\rangle\ \db\ \text{boy }\langle i+1,f-i-1\rangle
   \qquad(i=k,k+1,\ldots,f-1),
\]
ending at the girl~$\langle f,1\rangle$, who is free. If instead $f\le L$, then $k=1$,
and the search first walks backwards along hook~1 (each boy~$\langle 1,y\rangle$ with
$1+y>f$ has only his own girl) until it reaches the boy~$\langle 1,f-1\rangle$; from
there it descends the same staircase. In both cases the SAP ends at the
free girl of hook~$f$, the last girl in \A{QUEUE}, so step~\Hs{16} leaves the
order of the others unchanged.

\emph{The other roots fail.} The SAP passes through the girl~$\langle i,f+1-i\rangle$
of every hook~$i$ with $k<i\le f$. The root $\langle i,L\rangle$ of such a hook has only
one arc at each position~$y$ with $i+y>f$, namely to its own girl. Following
these single arcs, it reaches the boy~$\langle i,f+1-i\rangle$, whose only girl has been
deleted. Every list on the way is now empty, and step~\Hs{15} backs up all
the way: the root fails. So the round has exactly one SAP.

\emph{The new matching.} Augmenting dissolves, in each hook~$i$ with
$k<i<f$, the couple between positions $f-i$ and $f+1-i$, and couples
instead the boy at position~$f-i$ of hook~$i$ to the girl at position
$f-i$ of hook~$i+1$, for $k\le i<f$. Reading along the new couples, the
new hook~$i$ therefore consists of the first $f-i$ cells of the old hook
$i$, followed by the cells of the old hook~$i+1$ from position~$f-i$ on;
the remaining cells of the old hook~$i$ have passed to hook~$i-1$, and
the old hook~$f$ has been absorbed into hook~$f-1$. When $f\le L$, the
cells of hook~1 at positions $f,\ldots,L$, which the search walked
through, become coupled to themselves, that is, retired. Now compare
with the definition. In state~$f$, hook~$i+1$ has corner column~$f-i$, so
the new hook~$i$ runs east along row~$i$ to column~$f-i$ and then south:
it is hook~$i$ of state~$f-1$. The cells that retire are the alive cells
of column~$f$ other than the corner (which stays alive when $f=2$). Hooks~$1,\ldots,k-1$ are untouched whole rows, and they
are whole rows in state~$f-1$ too: for them $\ell_x\ge L+1=n+2-f$, so their
new corner column satisfies $f-x\ge n+2-x-\ell_x\ge\ell_x$, and if equality
holds then $x+2\ell_x=n+2$ and there is no cell below $(x,\ell_x)$. So the
new matching is $M_{f-1}$. For $f=1$ the search walks back along the single
hook and produces the perfect matching~$M_0$.
\end{proof}

Searched on its own, the root of each of the hooks~$k+1,\ldots,f$ would
find a SAP along its own hook. But \A{MARKED} lists the roots in hook
order, so the root of hook~$k$ goes first, and because the first arc of
each \A{BLINK} list points one hook down, its path takes the girl at
position~$f+1-i$ from every later hook~$i$. Each later root must pass
through that girl, so it fails.

\begin{lthm}{S}
On the ordered graph $\Gam_n$, Algorithm~H performs exactly $n$ rounds.
For $L=2,3,\ldots,n$, the round at level~$L$ starts with $f=n+1-L$ free
girls and finds exactly one SAP. The final matching consists of the
inside edges.
\end{lthm}

\begin{proof}
By Proposition~\ref{prop:gamma} and Lemma~\ref{lem:hooks}(d), step~\Hs1
produces $M_{n-1}$ with the free girls in hook order. Apply
Lemmas~\ref{lem:levels}--\ref{lem:dfs} for $f=n-1,n-2,\ldots,1$.
\end{proof}

\paragraph{Every size.}
Theorem~S gives $n$ rounds when $s=Q_n$. For $Q_n<s<Q_{n+1}$, append
$p=s-Q_n$ new girl--boy pairs to $\Gam_n$, numbered after all the old
ones. Give each new girl the list ``her own new boy, then the previous
boy,'' where the previous boy is the corner boy~$b_{n-1,2}$ for the first
new girl and the preceding new boy afterwards. Step~\Hs1 scans the new girls
first and couples each to her own boy. No old girl's list mentions a new
boy, so an alternating path from an old free girl never enters the new
pairs: to do so it would have to arrive at a boy whose mate is a new
girl. Hence every round runs exactly as on $\Gam_n$. The corner boy had
degree one, so the new graph is still planar, connected, and of maximum
degree three. For $s=1$ use a single edge, and for $s=0$ the empty
graph. Together with Section~\ref{sec:upper}, this proves Theorem~R.
\qed

\subsection{How much one path blocks}

In state~$f$ the hooks~$k,\ldots,f$ are $m=\min(f,L+1)$ disjoint SAPs, so
the largest set of disjoint SAPs has $\nu=m$ members, while Algorithm~H
takes one. The total number of SAPs in the dag grows exponentially in~$m$.

\begin{proposition}
In state~$f$, the dag of $\Gam_n$ contains exactly $2^m-1$ SAPs.
\end{proposition}

\begin{proof}
Let $N(x,y)$ be the number of dag paths from the boy at position~$y$ of
hook~$x$ down to $\bot$. By Lemma~\ref{lem:bfs},
$N(x,y)=N(x,y-1)+\Iv{x+y\le f}\,N(x+1,y-1)$, where $N(x,0)=1$ stands
for the free girl. Induction on $y$ gives $N(x,y)=2^{\min(y,f-x)}$. A root
$\langle i,L\rangle$ with $k\le i\le f$ therefore starts $2^{f-i}$ SAPs, since $f-i\le
f-k\le L$, and $\sum_{i=k}^{f}2^{f-i}=2^m-1$.
\end{proof}

Every SAP from the root of hook~$i>k$ passes through the girl at position
$f+1-i$ of that hook, which the first SAP uses. So the one SAP found by
Algorithm~H meets all $2^m-1$ SAPs, including all $m$ disjoint ones.
Lemma~\ref{lem:15ab}(b) with $\nu=m\le L+1$ in place of $f$ predicts at
least one SAP, and one is what Algorithm~H finds.

\subsection{Counting the work}
\label{sec:work}

\begin{lthm}{W}
In the round of $\Gam_n$ with $f$ free girls and level~$L=n+1-f$, with
$m=\min(f,L+1)$ as above,
\[
   W_5=3fL-n-\min(f,L),\qquad W_{13}=mL-\binom{m-1}{2},\qquad W_{11}=fL.
\]
Summed over all rounds after~\Hs1,
\[
  \sum W_5=\frac{(n-1)n(n+2)}2-Q_n+1,\qquad \sum W_{11}=\frac{(n-1)n(n+4)}6,
\]
\[
  \sum W_{13}=\begin{cases}(4u^3+21u^2+23u-6)/6,& n=2u+1,\\
                 (4u^3+15u^2+5u-12)/6,& n=2u.\end{cases}
\]
The final, unsuccessful search examines no edges.
\end{lthm}

\begin{proof}
The breadth-first search scans the girls at positions $1,\ldots,L$ of all
$f$ hooks, $fL$ girls in all, and every list entry of each. A girl has
three neighbours except that girls in column~1 have no left boy, girls in
row~1 have no upper boy, and the corner girl lacks the omitted edge. The
scanned girls include all $n-1$ girls of column~1, exactly $\min(f,L)$
girls of row~1, and the corner girl; this gives $W_5$. The dag contains
all $fL$ boys at positions $\le L$, by the last paragraph of the proof of
Lemma~\ref{lem:bfs}, so $W_{11}=fL$. The first root consumes $L$ arcs. The
root of hook~$k+h$, $1\le h<m$, consumes the arcs at positions
$L,L-1,\ldots,m-h$, including the one to a deleted girl: $L-m+h+1$
arcs. Adding up gives $W_{13}$, and summing over $f=1,\ldots,n-1$ gives the
totals (the sum $\sum_f\min(f,n+1-f)$ equals $Q_n-1$).
\end{proof}

For $\Gam_5$ the totals are $\sum W_5=62$, $\sum W_{13}=26$, $\sum W_{11}=30$, in agreement with the
table of Section~\ref{sec:watch}. As $n\to\infty$ we have $\sum W_5\sim n^3/2$,
$E\sim3n^2/4$, and $V=2Q_n\sim n^2/2$, so
\[
   \frac{\sum W_5}{E\sqrt V}\ \longrightarrow\ \frac{2\sqrt2}{3}\approx0.943 .
\]
Thus the breadth-first searches alone examine $\Theta(E\sqrt V)$
adjacency entries, and Corollary~K is tight up to a constant factor,
even for connected planar graphs of maximum degree three. The padding
used for other sizes adds only $O(n)$ edges and does not change this.
